\section*{Глава 2. Проектирование веб-приложения}
\addcontentsline{toc}{section}{Глава 2. Проектирование веб-приложения}
\stepcounter{section}


\subsection{Анализ целевой аудитории}

Разрабатываемое веб-приложение позиционируется как обучающий
симулятор (\emph{serious game}), предназначенный для подготовки
новых сотрудников фермерских хозяйств, специализирующихся на
выращивании личи (\textit{Litchi chinensis}). В отличие от
традиционных систем мониторинга, ориентированных на эксплуатацию
реального оборудования, данный продукт позволяет стажёрам
отрабатывать навыки управления поливом в безопасной виртуальной
среде, где ошибки не приводят к потере урожая и финансовым
потерям. Выделяются две ключевые роли пользователей:

\begin{itemize}
	\item \textbf{Агроном-наставник} --- опытный специалист,
	ответственный за создание учебных сценариев, конфигурацию
	секторов плантации, назначение стажёров на секторы, анализ
	журналов действий и оценку результатов обучения через таблицу
	лидеров и агрегированную телеметрию.
	\item \textbf{Стажёр-оператор полива} --- новый сотрудник,
	взаимодействующий с геймифицированным интерфейсом
	(виртуальная лейка, индексы здоровья, прогресс-бары,
	уведомления о критических событиях) для принятия решений
	о поливе на закреплённом секторе.
\end{itemize}

Для каждой роли проведено профилирование по методу персонажей
(Personas): агроном-наставник обладает высшим агрономическим
образованием, владеет понятийным аппаратом эвапотранспирации
и водного баланса, ожидает от системы физически осмысленную
динамику; стажёр-оператор имеет среднее образование, минимальный
опыт работы с цифровыми инструментами и нуждается в наглядной
обратной связи с элементами вовлечения (очки, бейджи, анимации).
Исследование Shen et~al.~(2024)~\cite{shen2024_gamification_workforce}
экспериментально подтвердило, что комбинация аватаров, очков,
бейджей, мгновенной обратной связи и таблиц лидеров значимо
повышает самоэффективность и закрепление навыков у работников
в корпоративном обучении.


\subsection{Описание функциональности приложения}

Функциональные возможности системы определяются через взаимодействие
пользователя с симулятором. Основной обучающий сценарий представлен
в виде \textit{User Story}: <<Я, как стажёр-оператор, хочу видеть
текущий уровень влажности почвы, температуру, фенофазу растения и
индекс здоровья на моём секторе, чтобы принять решение о поливе
через интерфейс виртуальной лейки и получить мгновенную обратную
связь о правильности своих действий>>.

На рисунке~\ref{fig:use_case} представлена диаграмма вариантов
использования (Use Case).

\begin{figure}[ht!]
	\centering
	\includegraphics[width=0.9\textwidth]{UML}
	\caption{Диаграмма вариантов использования системы}
	\label{fig:use_case}
\end{figure}

Ключевые функциональные блоки:
\begin{enumerate}
	\item \textbf{Управление секторами и сценариями} (агроном):
	создание секторов, задание начальных параметров почвы и
	климата, выбор фенофазы личи, настройка вероятностей
	случайных событий, импорт конфигурации из CSV.
	\item \textbf{Интерактивный полив} (стажёр): выбор объёма
	воды и момента полива, визуализация последствий через
	анимацию лейки и динамику прогресс-баров.
	\item \textbf{Мониторинг в реальном времени}: потоковая
	передача состояния симуляции по WebSocket, графики
	временных рядов телеметрии.
	\item \textbf{Обучающая аналитика}: журнал действий стажёра,
	сравнение с эталонным сценарием полива, таблица лидеров,
	система бейджей.
	\item \textbf{Экспорт данных}: выгрузка телеметрии в XLSX
	для офлайн-анализа наставником.
\end{enumerate}


\subsection{Проектирование математической модели симуляции}

Центральным элементом проектирования является серверный движок
симуляции, реализующий физически обоснованную математическую
модель агрономических процессов. Движок работает в тиковом
режиме: каждый <<тик>> соответствует одному игровому дню
(настраиваемый интервал реального времени --- от 10~секунд до
нескольких минут). Модель включает шесть взаимосвязанных
подсистем, последовательно выполняемых на каждом тике.

\subsubsection{Подсистема генерации погоды}

Для создания реалистичных и неповторяющихся обучающих сценариев
погодные условия генерируются стохастически на основе цепи Маркова
первого порядка~\cite{richardson1981_weather,semenov1998_larswg}.
Состояние <<сухой/влажный день>> определяется матрицей переходных
вероятностей:

\begin{equation}
	P = \begin{pmatrix} P_{00} & P_{01} \\ P_{10} & P_{11} \end{pmatrix},
	\label{eq:markov_weather}
\end{equation}

\begin{where}
	\item $P_{01}$~--- вероятность перехода из сухого дня во влажный;
	\item $P_{11}$~--- вероятность продолжения влажного дня;
	\item $P_{00} = 1 - P_{01}$, $P_{10} = 1 - P_{11}$.
\end{where}

Объём осадков во влажный день моделируется через гамма-распределение
$\Gamma(\alpha, \beta)$ с параметрами формы $\alpha$ и масштаба
$\beta$, откалиброванными по климатическим данным региона
(Гватемала, тропический пояс). Суточные значения максимальной
и минимальной температуры ($T_{max}$, $T_{min}$) генерируются
авторегрессионной моделью первого порядка, условно на состояние
<<мокрый/сухой>>~\cite{chen2012_weagets,chen2020_markov_bias}.

Данный подход обеспечивает вариативность обучающих сессий: при
каждом новом прохождении стажёр сталкивается с уникальной
последовательностью погодных условий, что исключает заучивание
фиксированного сценария и формирует навык адаптивного принятия
решений.

\subsubsection{Подсистема расчёта эвапотранспирации}

Суточная потеря воды из почвы рассчитывается на основе эталонной
эвапотранспирации $ET_0$. В базовой конфигурации движка используется
упрощённое уравнение Харгривса--Самани~\cite{hargreaves2003_history},
требующее только температурных
данных~\cite{aschonitis2017_modified_hargreaves}:

\begin{equation}
	ET_{0} = 0{,}0023 \cdot R_a \cdot (T_{mean} + 17{,}8) \cdot (T_{max} - T_{min})^{0{,}5},
	\label{eq:hargreaves}
\end{equation}

\begin{where}
	\item $R_a$~--- внеатмосферная радиация, МДж$\cdot$м$^{-2}\cdot$сут$^{-1}$ (рассчитывается по широте и дню года);
	\item $T_{mean} = (T_{max} + T_{min}) / 2$~--- среднесуточная температура воздуха, °C;
	\item $T_{max}$, $T_{min}$~--- суточные экстремумы температуры, °C.
\end{where}

Для расширенной конфигурации предусмотрена реализация полного
уравнения Пенмана--Монтейса (FAO-56)~\cite{allen1998_fao56,
	zotarelli2010_step_by_step}:

\begin{equation}
	ET_0 = \frac{0{,}408 \cdot \Delta \cdot (R_n - G) + \gamma \cdot \dfrac{900}{T + 273} \cdot u_2 \cdot (e_s - e_a)}{\Delta + \gamma \cdot (1 + 0{,}34 \cdot u_2)},
	\label{eq:penman_monteith}
\end{equation}

\begin{where}
	\item $\Delta$~--- наклон кривой давления насыщенного водяного пара, кПа$\cdot$°C$^{-1}$;
	\item $R_n$~--- чистая радиация на поверхности, МДж$\cdot$м$^{-2}\cdot$сут$^{-1}$;
	\item $G$~--- поток тепла в почву, МДж$\cdot$м$^{-2}\cdot$сут$^{-1}$;
	\item $\gamma$~--- психрометрическая постоянная, кПа$\cdot$°C$^{-1}$;
	\item $u_2$~--- скорость ветра на высоте 2~м, м/с;
	\item $e_s - e_a$~--- дефицит давления насыщения пара, кПа.
\end{where}

Фактическая эвапотранспирация культуры определяется через дуальный
коэффициент~\cite{allen2005_dual_kc}:

\begin{equation}
	ET_c = (K_s \cdot K_{cb} + K_e) \cdot ET_0,
	\label{eq:dual_kc}
\end{equation}

\begin{where}
	\item $K_{cb}$~--- базальный коэффициент культуры (зависит от текущей фенофазы личи);
	\item $K_e$~--- коэффициент испарения с поверхности почвы;
	\item $K_s$~--- коэффициент водного стресса ($0 \leq K_s \leq 1$).
\end{where}

\subsubsection{Подсистема водного баланса почвы}

Состояние влажности почвы описывается через дефицит влаги в
корневой зоне $D_r$, обновляемый на каждом тике по уравнению
водного баланса~\cite{allen1998_fao56,thorp2022_pyfao56}.
Схема потоков воды в системе <<атмосфера--почва--корневая
зона--дренаж>> представлена на рисунке~\ref{fig:water_balance}.

\begin{figure}[ht!]
	\centering
	\includegraphics[width=0.55\textwidth]{water_balance_scheme}
	\caption{Схема водного баланса почвы в модели симуляции}
	\label{fig:water_balance}
\end{figure}

Уравнение водного баланса:

\begin{equation}
	D_r(t+1) = D_r(t) - P_{eff}(t) - I(t) + ET_c(t) + DP(t),
	\label{eq:water_balance}
\end{equation}

\begin{where}
	\item $D_r(t)$~--- дефицит влаги на начало тика, мм;
	\item $P_{eff}(t)$~--- эффективные осадки (за вычетом поверхностного стока), мм;
	\item $I(t)$~--- объём полива, внесённый оператором, мм;
	\item $ET_c(t)$~--- фактическая эвапотранспирация, мм;
	\item $DP(t)$~--- глубинная перколяция (дренаж), мм.
\end{where}

Коэффициент водного стресса $K_s$ вычисляется по методологии
FAO-56 через сопоставление текущего дефицита с общим доступным
запасом воды (\emph{Total Available Water}, TAW) и пороговой
долей истощения $p$~\cite{raes2023_aquacrop_ch3,
	steduto2012_aquacrop_ch1}:

\begin{equation}
	K_s = \begin{cases}
			  1, & \text{если } D_r \leq p \cdot TAW, \\[6pt]
			  \dfrac{TAW - D_r}{(1 - p) \cdot TAW}, & \text{если } D_r > p \cdot TAW.
	\end{cases}
	\label{eq:ks_stress}
\end{equation}

\indent Значение $K_s = 1$ означает отсутствие водного стресса; при
$K_s \to 0$ растение испытывает критическую засуху. Параметр
$p$ зависит от культуры и фенофазы (для тропических плодовых
$p \approx 0{,}5$).

\subsubsection{Подсистема фенологии растений}

Развитие личи моделируется через накопление суммы эффективных
температур (\emph{Growing Degree Days},
GDD)~\cite{wikipedia_gdd,bygl2024_gdd_phenology}:

\begin{equation}
	GDD(t) = \max\!\left(0, \;\frac{T_{max}(t) + T_{min}(t)}{2} - T_{base}\right),
	\label{eq:gdd}
\end{equation}

\begin{where}
	\item $T_{base}$~--- базовая температура развития ($\approx 10$~°C для тропических плодовых);
	\item $T_{max}(t)$, $T_{min}(t)$~--- суточные экстремумы температуры.
\end{where}

Кумулятивная сумма $\sum GDD$ определяет переключение между
фенофазами личи по расширенной шкале
BBCH~\cite{wei2013_litchi_bbch}: разрыв почек $\to$ развитие
побегов $\to$ появление соцветий $\to$ цветение $\to$ развитие
плодов $\to$ созревание. Kumar et~al.~(2025)~\cite{kumar2025_litchi_growth}
детализировали 74 вторичные стадии для культивара <<Dehradun>>
и сопоставили каждую стадию с накоплением GDD, что предоставляет
калибровочные данные для симулятора.

На каждой фенофазе в модели автоматически обновляются:
коэффициент $K_{cb}$ (базальный коэффициент культуры в
формуле~(\ref{eq:dual_kc})), чувствительность к водному
стрессу и оптимальный диапазон влажности. Это делает задачу
определения правильного момента и объёма полива нетривиальной,
что является основой обучающего сценария.

\subsubsection{Подсистема индекса здоровья растений}

Индекс здоровья $H \in [0; 1]$ интегрирует влияние всех
стрессовых факторов и является центральным показателем,
визуализируемым в геймифицированном интерфейсе. Обновление
индекса на каждом тике выполняется по
формуле~\cite{shaw2015_waterlogging,liu2020_waterlogging_review,
	raes2023_aquacrop_ch3}:

\begin{equation}
	\begin{split}
		H(t+1) = \text{clamp}_{[0,\,1]}\!\Big(
		H(t) & - \alpha_{drought} \cdot (1 - K_s)
		- \alpha_{flood} \cdot (1 - K_{s,aer}) \\
		& - \alpha_{heat} \cdot (1 - K_{s,T})
		+ \beta_{rec} \cdot K_s \cdot K_{s,aer} \cdot (1 - H(t))
		\Big),
	\end{split}
	\label{eq:health_index}
\end{equation}

\begin{where}
	\item $\alpha_{drought}$~--- коэффициент деградации при засухе (калибруется для геймплея, $\approx 0{,}03$);
	\item $\alpha_{flood}$~--- коэффициент деградации при переувлажнении ($\approx 0{,}02$);
	\item $\alpha_{heat}$~--- коэффициент деградации при тепловом стрессе ($\approx 0{,}01$);
	\item $K_{s,aer}$~--- коэффициент аэрационного стресса (переувлажнение), рассчитываемый по аналогии с AquaCrop;
	\item $K_{s,T}$~--- температурный коэффициент стресса;
	\item $\beta_{rec}$~--- коэффициент восстановления ($\approx 0{,}05$), активный только при отсутствии стрессов ($K_s \approx 1$, $K_{s,aer} \approx 1$);
	\item $\text{clamp}_{[0,\,1]}$~--- ограничение значения диапазоном $[0; 1]$.
\end{where}

Динамика переходов между состояниями сектора (нормальное,
засуха, переувлажнение, тепловой стресс, критическое,
восстановление, гибель) представлена на
рисунке~\ref{fig:state_machine} в виде диаграммы состояний UML.

\begin{figure}[ht!]
	\centering
	\includegraphics[width=0.9\textwidth]{state_machine}
	\caption{Диаграмма состояний сектора плантации}
	\label{fig:state_machine}
\end{figure}

Коэффициент аэрационного стресса $K_{s,aer}$ описывает реакцию
корневой системы на переувлажнение. По методологии
AquaCrop~\cite{raes2023_aquacrop_ch3} и трёхэтапной модели
Shaw~\&~Meyer~(2015)~\cite{shaw2015_waterlogging}:

\begin{equation}
	K_{s,aer} = \begin{cases}
					1, & \text{если } \theta \leq \theta_{SAT} - \theta_{aer}, \\[6pt]
					\dfrac{\theta_{SAT} - \theta}{\theta_{aer}}, & \text{если } \theta > \theta_{SAT} - \theta_{aer},
	\end{cases}
	\label{eq:ks_aer}
\end{equation}

\begin{where}
	\item $\theta$~--- текущая объёмная влажность почвы;
	\item $\theta_{SAT}$~--- влажность насыщения;
	\item $\theta_{aer}$~--- порог аэрации (объём воздуха, необходимый корням).
\end{where}

Дополнительно в модели реализован индекс водного стресса
растений (\emph{Crop Water Stress Index},
CWSI)~\cite{jackson1988_cwsi,idso1981_cwsi}, используемый
для визуализации в интерфейсе агронома:

\begin{equation}
	CWSI = \frac{(T_c - T_{air}) - (T_c - T_{air})_{LL}}{(T_c - T_{air})_{UL} - (T_c - T_{air})_{LL}},
	\label{eq:cwsi}
\end{equation}

\begin{where}
	\item $T_c$~--- температура кроны (моделируется на основе $K_s$ и радиации);
	\item $(T_c - T_{air})_{LL}$~--- нижняя базовая линия (отсутствие стресса);
	\item $(T_c - T_{air})_{UL}$~--- верхняя базовая линия (максимальный стресс).
\end{where}

Значение $CWSI = 0$ соответствует полностью обводнённому
растению, $CWSI = 1$ --- критическому водному
стрессу~\cite{zhou2024_cwsi_maize,katimbo2021_cwsi_rootzone}.

\subsubsection{Подсистема случайных событий (Монте-Карло)}

Для формирования навыков реагирования на нештатные ситуации
в симуляцию интегрирован модуль стохастических событий,
реализованный методом Монте-Карло~\cite{sun2025_monte_carlo_crop,
	jones2017_bayesian_sugarcane,graveline2012_monte_carlo_farming}.
На каждом тике выполняется бернуллиевская выборка для событий:

\begin{equation}
	\Pr(\text{событие}_i) = p_i(\Sigma GDD, \; \text{фенофаза}, \; \text{сезон}),
	\label{eq:event_probability}
\end{equation}

\begin{where}
	\item $p_i$~--- вероятность наступления $i$-го события, зависящая от накопленных GDD и текущей фенофазы.
\end{where}

Реализованы следующие типы случайных событий:
\begin{itemize}
	\item \textbf{Нашествие вредителей} --- вероятность
	пропорциональна накопленным GDD от биологического
	порога~\cite{vieira2025_wheat_sawfly_gdd}; при наступлении
	снижается $H$ на $\Delta H_{pest}$, стажёр должен
	активировать <<обработку>>.
	\item \textbf{Аномальная жара} --- суточная вероятность
	$p_{heat} \approx 0{,}05$ в сухой сезон; повышает
	$T_{max}$ на 5--8~°C, увеличивая $ET_0$ и ускоряя
	истощение влаги.
	\item \textbf{Ливень} --- вероятность определяется цепью
	Маркова (формула~(\ref{eq:markov_weather})); при
	избыточных осадках активируется $K_{s,aer}$, стажёр
	должен воздержаться от полива.
	\item \textbf{Поломка оборудования} --- постоянная
	суточная вероятность $p_{equip} \approx 0{,}02$;
	ограничивает доступный объём полива на 1--3~тика.
\end{itemize}

Вариативность событий обеспечивает уникальность каждого
прохождения и формирует у стажёров навык адаптивного
принятия решений в условиях
неопределённости~\cite{newlands2014_probabilistic_yield}.


\subsection{Проектирование механик геймификации}

Геймификация в разрабатываемом приложении выполняет функцию
обучающего инструмента, а не развлекательного дополнения.
На основе систематического обзора серьёзных игр в сельском
хозяйстве~\cite{madeira2024_serious_games_agriculture,
	cinar2024_serious_games_scoping} и исследований геймификации
в корпоративном обучении~\cite{shen2024_gamification_workforce}
спроектированы следующие механики:

\begin{enumerate}
	\item \textbf{Система очков.} Начисляются за поддержание
	индекса здоровья $H > 0{,}7$, экономию воды (полив в
	пределах суточного лимита), своевременное реагирование
	на критические события. Штрафные очки --- за допущение
	засухи ($K_s < 0{,}3$) или переувлажнения
	($K_{s,aer} < 0{,}5$).
	\item \textbf{Система бейджей.} Присуждаются за
	достижение пороговых значений: <<Хранитель воды>>
	(7~дней без перелива), <<Зелёный мастер>> ($H > 0{,}9$
	в течение 14~дней), <<Кризис-менеджер>> (успешное
	преодоление 3~случайных событий подряд).
	\item \textbf{Прогресс-бары и визуальная обратная связь.}
	Индекс здоровья, уровень влажности почвы, суточный лимит
	воды и накопленные GDD отображаются в виде анимированных
	прогресс-баров с цветовой индикацией (зелёный $\to$
	жёлтый $\to$ красный).
	\item \textbf{Таблица лидеров.} Ранжирование стажёров
	по совокупному баллу за обучающую сессию; доступна
	агроному-наставнику для сравнительного анализа.
	\item \textbf{Немедленная обратная связь.} WebSocket-уведомления о критических событиях
	(засуха, переувлажнение, падение $H$ ниже порога,
	появление вредителей) с рекомендацией действий.
\end{enumerate}

Исследование Suwarno et~al.~(2014)~\cite{suwarno2014_virtual_farmers}
продемонстрировало эффективность комбинации исторической
симуляции и игрового повествования для обучения фермеров;
аналогичный подход реализован в данном проекте через
сценарные <<миссии>> с нарастающей сложностью.


\subsection{Проектирование архитектуры движка симуляции}

Серверный движок реализует тиковую архитектуру
(\emph{tick-based simulation engine}), канонический цикл
которой описан Ropert~(2025)~\cite{ropert2025_game_tick}:

\begin{enumerate}
	\item Приём входящих сообщений (действия оператора: полив, обработка от вредителей).
	\item Диспетчеризация команд к целевым сущностям (секторам).
	\item Эволюция мира на один тик (последовательное выполнение подсистем).
	\item Постановка обновлений в очередь рассылки клиентам через WebSocket.
\end{enumerate}

На рисунке~\ref{fig:tick_pipeline} представлена схема конвейера
подсистем, выполняемых за один тик симуляции.

\begin{figure}[ht!]
	\centering
	\includegraphics[width=\textwidth]{tick_pipeline}
	\caption{Конвейер подсистем движка симуляции (один тик)}
	\label{fig:tick_pipeline}
\end{figure}

Подсистемы выполняются в строгом порядке:
WeatherTick $\to$ ET0Tick $\to$ WaterBalanceTick $\to$
PhenologyTick $\to$ StressTick $\to$ HealthTick $\to$
EventTick $\to$ ScoringTick. Такая цепочка обеспечивает
причинно-следственную согласованность: погода определяет
эвапотранспирацию, эвапотранспирация влияет на водный баланс,
водный баланс определяет стресс, стресс обновляет здоровье,
а случайные события зависят от
фенофазы~\cite{getgud2024_multiplayer_arch,
	techtide2024_game_server}.

Каждый сектор плантации моделируется как независимая сущность
с собственным состоянием (профиль почвы, глубина корней, индекс
здоровья, фенофаза, накопленные GDD). Это позволяет параллельно
обрабатывать множество секторов в отдельных горутинах Go.
Детерминированность симуляции (при фиксированном seed генератора
случайных чисел) обеспечивает воспроизводимость обучающих
сценариев для аналитики действий стажёра.


\subsection{Проектирование модели данных}

Структура базы данных спроектирована с учётом необходимости хранения
конфигураций секторов, журналов полива, временных рядов телеметрии,
а также обучающей аналитики (очки, бейджи, результаты сессий).
В качестве СУБД используется PostgreSQL~16.

На рисунке~\ref{fig:er_diagram} представлена ER-диаграмма базы
данных.

\begin{figure}[ht!]
	\centering
	\includegraphics[width=\textwidth]{er_diagram}
	\caption{Диаграмма <<сущность--связь>> базы данных}
	\label{fig:er_diagram}
\end{figure}

Структура таблицы секторов плантации представлена в табл.~\ref{tab:db_sectors}.

\begin{longtable}{|l|l|p{3.5cm}|p{5.5cm}|}
	\caption{Структура таблицы секторов (\texttt{sectors})} \label{tab:db_sectors}           \\

	\hline
	\textbf{Поле}       & \textbf{Тип}      & \textbf{Ограничение}    & \textbf{Описание}              \\ \hline
	\endfirsthead

	\tabcont{\ref{tab:db_sectors}}
	\textbf{Поле}       & \textbf{Тип}      & \textbf{Ограничение}    & \textbf{Описание}              \\ \hline
	\endhead

	\hline
	\endfoot

	\hline
	\endlastfoot

	id                   & UUID              & PRIMARY KEY             & Уникальный идентификатор сектора   \\ \hline
	name                 & VARCHAR(255)      & NOT NULL                & Название сектора                   \\ \hline
	area\_sqm            & FLOAT             & DEFAULT 0               & Площадь сектора, м$^2$             \\ \hline
	soil\_moisture       & FLOAT             & DEFAULT 50              & Текущая влажность почвы, \%        \\ \hline
	temperature          & FLOAT             & DEFAULT 25              & Температура воздуха, °C            \\ \hline
	health\_index        & FLOAT             & DEFAULT 1.0             & Индекс здоровья $H$ (0--1)         \\ \hline
	gdd\_cumulative      & FLOAT             & DEFAULT 0               & Накопленные GDD                    \\ \hline
	phenophase           & VARCHAR(10)       & DEFAULT '00'            & Код фенофазы BBCH                  \\ \hline
	ks\_water            & FLOAT             & DEFAULT 1.0             & Коэффициент $K_s$ (водный стресс)  \\ \hline
	ks\_aeration         & FLOAT             & DEFAULT 1.0             & Коэффициент $K_{s,aer}$ (аэрация)  \\ \hline
	deficit\_dr          & FLOAT             & DEFAULT 0               & Дефицит влаги $D_r$, мм            \\ \hline
	status               & VARCHAR(50)       & DEFAULT 'normal'        & Статус (normal, drought, overwatered) \\ \hline
	operator\_id         & UUID              & FK $\to$ users, NULL    & Назначенный стажёр                 \\ \hline
	last\_watered\_at    & TIMESTAMPTZ       & NULL                    & Время последнего полива            \\ \hline
	daily\_water\_limit  & FLOAT             & DEFAULT 500             & Лимит воды на сутки, л             \\ \hline
	water\_consumed      & FLOAT             & DEFAULT 0               & Израсходовано воды за сутки, л     \\ \hline
	created\_at          & TIMESTAMPTZ       & DEFAULT NOW()           & Время создания записи              \\ \hline
	updated\_at          & TIMESTAMPTZ       & DEFAULT NOW()           & Время последнего изменения         \\ \hline
\end{longtable}

Помимо таблицы секторов, база данных включает таблицы пользователей
(\texttt{users}) с ролевой моделью (\texttt{agronomist},
\texttt{operator}), растений (\texttt{plants}), журнала полива
(\texttt{watering\_logs}), телеметрии (\texttt{telemetry}),
уведомлений (\texttt{notifications}), а также таблицу обучающей
аналитики (\texttt{training\_scores}) с полями: \texttt{user\_id},
\texttt{session\_id}, \texttt{total\_score}, \texttt{badges}
(массив JSONB), \texttt{avg\_health}, \texttt{water\_efficiency}.

Таблица телеметрии индексирована составным индексом по полям
\texttt{sector\_id} и \texttt{recorded\_at~DESC} для эффективной
выборки временных рядов. Параметры погодного генератора
(матрица переходов, параметры гамма-распределения) хранятся в
таблице \texttt{weather\_configs} с возможностью задания профилей
для различных климатических зон.


\subsection{Разработка макетов страниц}

На основе анализа целевой аудитории и функциональных требований
разработаны макеты (\emph{wireframes}) ключевых экранов
приложения в формате low-fidelity.

\subsubsection{Дашборд стажёра-оператора}

На рисунке~\ref{fig:wireframe_operator} представлен макет
главного экрана стажёра-оператора. Интерфейс включает четыре
карточки-индикатора (влажность почвы, температура воздуха,
индекс здоровья, суточный лимит воды), кнопку полива с
регулятором объёма, мини-график телеметрии за последние
24~часа, ленту уведомлений и список полученных бейджей.

\begin{figure}[ht!]
	\centering
	\includegraphics[width=\textwidth]{wireframe_operator}
	\caption{Макет дашборда стажёра-оператора}
	\label{fig:wireframe_operator}
\end{figure}

\subsubsection{Дашборд агронома-наставника}

Макет интерфейса агронома (рисунок~\ref{fig:wireframe_agronomist})
ориентирован на мониторинг множества секторов одновременно.
Основная область содержит таблицу секторов с ключевыми
показателями (здоровье, влажность, статус, фенофаза), кнопки
управления (создание сектора, импорт CSV, экспорт XLSX),
таблицу лидеров и сводный график среднего индекса здоровья
по всем секторам.

\begin{figure}[ht!]
	\centering
	\includegraphics[width=\textwidth]{wireframe_agronomist}
	\caption{Макет дашборда агронома-наставника}
	\label{fig:wireframe_agronomist}
\end{figure}

\subsubsection{Экран управления поливом}

Развёрнутый экран полива (рисунок~\ref{fig:wireframe_watering})
представляет полноэкранную визуализацию сектора в виде
схематичной плантации, где цвет деревьев зависит от текущего
индекса здоровья $H$ (зелёный --- норма, жёлтый --- стресс,
коричневый --- критическое состояние). Анимированная лейка
в центре активируется при нажатии кнопки <<Полить>>. Боковая
панель отображает текущие значения коэффициентов $K_s$,
$K_{s,aer}$, $CWSI$ и $GDD$ с всплывающими подсказками,
поясняющими формулы расчёта. Всплывающие уведомления
(\emph{toast}) информируют о критических событиях.

\begin{figure}[ht!]
	\centering
	\includegraphics[width=\textwidth]{wireframe_watering}
	\caption{Макет экрана управления поливом}
	\label{fig:wireframe_watering}
\end{figure}